\documentclass[11pt]{article}
\usepackage[margin=1in]{geometry}
\usepackage{amsmath,amssymb,amsthm}
\usepackage[colorlinks=true,linkcolor=blue,urlcolor=blue,citecolor=blue]{hyperref}
\newtheorem{theorem}{Theorem}[section]
\newtheorem{lemma}[theorem]{Lemma}
\newtheorem{proposition}[theorem]{Proposition}
\newtheorem{remark}[theorem]{Remark}
\newcommand{\tr}{\operatorname{tr}}
\newcommand{\diag}{\operatorname{diag}}
\newcommand{\N}{\mathcal N}
\newcommand{\one}{\boldsymbol 1}
\newcommand{\R}{\mathbb R}
\title{The local and global determinant budget for a doubled unit edge:\\
sharp coefficients at every fixed maximum degree}
\author{Research session supplement; frozen manuscripts unchanged}
\date{November 2026}
\begin{document}
\maketitle
\begin{abstract}
For the relative fractional-Laplacian correction produced by doubling
one unit bridge in a finite simple unweighted graph of maximum degree
three, we derive the worst-case count asymptotic
$\N(\varepsilon)\sim\alpha\log(6/5)/\varepsilon$ after taking the
supremum over hosts. More generally the coefficient at maximum degree
$d\ge3$ is $\alpha\log(2d/(d+2))$, also for bipartite hosts.
A free/wired resistance comparison extends the same sharp coefficient
from bridge edits to arbitrary existing-edge doubling.
The upper bound charges the entire rooted local
response against the finite determinant budget. Its proof includes an
ordinary trace-class Fredholm identity without a spectral-gap assumption
and a generalized Ritz transfer involving only positive fractional
powers. For the lower bound, many ports of each reservoir share one
mass coordinate. Repeated binary-tree blocks dilute their contact
resistances and realize a tunable weighted star. Its reverse pivots
allocate any amount strictly below $\log(6/5)$ equally among arbitrarily
many modes. All examples remain finite, simple, and unweighted before the
update. This is a new research proof,
not an edit to or a claim of external verification of the frozen papers.
\end{abstract}

\section{Statement and scope}
Fix $0<\alpha<1$. For a finite connected simple unit-weight graph $G$
and an existing bridge $e=\{u,v\}$, write $A=L(G)$, $b=e_u-e_v$,
$B=A+bb^*$, and
\[
 C_\alpha(G,e)=(A^\dagger)^{\alpha/2}(B^\alpha-A^\alpha)
                                      (A^\dagger)^{\alpha/2}.
\]
Pseudoinverse powers vanish on constants. Set
\[
 \N_C(\varepsilon)=\#\{j:\lambda_j(C)>\varepsilon\},\qquad
 W^{\mathrm{br}}_{d,\alpha}(\varepsilon)
   =\sup_{\substack{\deg G\le d\\e\text{ a bridge}}}
                                    \N_{C_\alpha(G,e)}(\varepsilon).
\]
Combinatorial degree is measured before the update.
Write $W^{\mathrm{all}}_{d,\alpha}$ for the same supremum allowing
any existing edge to be doubled.

\begin{theorem}\label{thm:sharp}
For every fixed $0<\alpha<1$ and integer $d\ge3$,
\[
 \lim_{\varepsilon\downarrow0}\varepsilon
               W^{\mathrm{br}}_{d,\alpha}(\varepsilon)
 =\lim_{\varepsilon\downarrow0}\varepsilon
               W^{\mathrm{all}}_{d,\alpha}(\varepsilon)
 =\alpha\log\frac{2d}{d+2}.
\]
The same equality holds when the supremum is restricted to bipartite
hosts. In particular the cubic coefficient is $\alpha\log(6/5)$.
\end{theorem}

The bridge half-tree calculation alone does not apply to nonbridges;
Section~\ref{sec:free-wired} supplies the additional comparison needed
for that extension. No efficiency or moderate-size claim is made for
the existential lower examples.

\section{The exact local Fredholm determinant, including gap zero}
\begin{lemma}\label{lem:fredholm}
Let $A\ge0$ be bounded, self-adjoint and injective on a separable Hilbert
space. Suppose $\tau=\|A^{-1/2}b\|^2<\infty$, and put $B=A+bb^*$.
There is a unique bounded positive operator $C$ satisfying
\[
 B^\alpha-A^\alpha=A^{\alpha/2}CA^{\alpha/2}.
\]
It is trace class, $\tr C\le\alpha\tau$, and
\begin{equation}\label{eq:fredholm}
 \log\det_F(I+C)=\tr\log(I+C)=\alpha\log(1+\tau).
\end{equation}
No lower spectral bound on $A$ is required.
\end{lemma}
\begin{proof}
Let $c_\alpha=\sin(\pi\alpha)/\pi$ and, for $t>0$, define
\[
 m(t)=\langle b,(A+t)^{-1}b\rangle,\qquad
 w_t=A^{-\alpha/2}(A+t)^{-1}b.
\]
These vectors exist: in the spectral measure $\mu$ of $b$,
$\lambda^{-\alpha}$ is bounded by a constant times $\lambda^{-1}$
on the bounded positive spectrum. Tonelli and the scalar beta integral
give
\begin{equation}\label{eq:trace-majorant}
 c_\alpha\int_0^\infty t^\alpha\|w_t\|^2\,dt
       =\alpha\int\lambda^{-1}\,d\mu(\lambda)=\alpha\tau.
\end{equation}
Thus
\begin{equation}\label{eq:resolvent}
 C=c_\alpha\int_0^\infty\frac{t^\alpha w_tw_t^*}{1+m(t)}\,dt
\end{equation}
is a positive trace-norm Bochner integral. Multiplying by
$A^{\alpha/2}$ on both sides and using the fractional resolvent formula
and the rank-one resolvent identity proves the stated fractional energy
identity. The range of $A^{\alpha/2}$ is dense, proving uniqueness.

For the determinant, regularize $A_\delta=A+\delta I$, $\delta>0$,
and form $C_\delta$ for the pair $A_\delta,A_\delta+bb^*$.
In~\eqref{eq:resolvent} replace $w_t$ by
\[
 w_{\delta,t}=(A+\delta)^{-\alpha/2}(A+\delta+t)^{-1}b.
\]
For each $t>0$, $w_{\delta,t}\to w_t$ and
$\|w_{\delta,t}\|\le\|w_t\|$; the scalar denominators converge too.
The rank-one trace norm inequality
$\|xx^*-yy^*\|_1\le(\|x\|+\|y\|)\|x-y\|$ and
\eqref{eq:trace-majorant} imply
\begin{equation}\label{eq:s1-convergence}
                     \|C_\delta-C\|_1\longrightarrow0.
\end{equation}
For fixed $\delta$, set
\[
 X_s=A_\delta+sbb^*,\qquad
 R_s=A_\delta^{-\alpha/2}X_s^\alpha A_\delta^{-\alpha/2},
                  \quad 0\le s\le1.
\]
This is an invertible identity-plus-trace-class $C^1$ family. Its
derivative follows by differentiating the fractional resolvent formula:
\[
 \dot X_s^\alpha=c_\alpha\int_0^\infty
        t^\alpha(X_s+t)^{-1}bb^*(X_s+t)^{-1}\,dt.
\]
The common lower bound $\delta$ justifies all trace-norm derivatives
and integrals. Fredholm logarithmic differentiation and trace cyclicity
give
\begin{align*}
 \frac{d}{ds}\log\det_F R_s
 &=\tr\bigl(X_s^{-\alpha}\dot X_s^\alpha\bigr)\\
 &=c_\alpha\int_0^\infty t^\alpha
        \langle b,X_s^{-\alpha}(X_s+t)^{-2}b\rangle\,dt\\
 &=\alpha\langle b,X_s^{-1}b\rangle
   =\frac{\alpha\tau_\delta}{1+s\tau_\delta},
       \qquad\tau_\delta=\langle b,A_\delta^{-1}b\rangle.
\end{align*}
Integration yields $\log\det_F(I+C_\delta)
=\alpha\log(1+\tau_\delta)$. Now $\tau_\delta\uparrow\tau$, and
trace-norm continuity of the Fredholm determinant together with
\eqref{eq:s1-convergence} proves~\eqref{eq:fredholm}.
All negative powers in the unregularized proof act on declared domain
vectors; the determinant computation itself uses bounded inverses.
Standard determinant definitions and continuity are recalled in
\cite[Sections 3--4]{Bornemann}.
\end{proof}

\section{Rooted convergence forces the whole local budget}
\begin{proposition}\label{prop:ritz}
Suppose finite edge-rooted graphs $(G_n,e_n)$ of uniformly bounded
degree converge locally to an infinite connected $(G,e)$. Let $C_n$
be their relative corrections and $C$ the operator of
Lemma~\ref{lem:fredholm} for $A=L(G)$, $b=e_u-e_v$. Then
\[
                 \lambda_j(C)\le\liminf_n\lambda_j(C_n)
                 \quad\text{for every fixed }j.
\]
\end{proposition}
\begin{proof}
The infinite connected graph Laplacian is injective on $\ell^2$:
zero energy forces a constant function, and an $\ell^2$ constant on
an infinite graph is zero. The edge energy gives $bb^*\preceq A$,
so its wired leverage $\tau=\|A^{-1/2}b\|^2\le1$ and
Lemma~\ref{lem:fredholm} applies.

Take linearly independent finitely supported $f_1,\ldots,f_k$ and
transport them into the identified rooted balls of $G_n$. Polynomial
locality and uniform polynomial approximation to $x^\alpha$ on a
common bounded spectral interval imply convergence of both Gram
matrices
\[
 E_n(i,j)=\langle f_{i,n},A_n^\alpha f_{j,n}\rangle,
 \quad D_n(i,j)=\langle f_{i,n},(B_n^\alpha-A_n^\alpha)f_{j,n}\rangle
\]
to $E,D$ on $G$. Since $A$ is injective, $E\succ0$, and hence
$E_n\succ0$ eventually. Therefore
$E_n^{-1/2}D_nE_n^{-1/2}\to E^{-1/2}DE^{-1/2}$.
These are precisely the Ritz matrices on the images under
$A_n^{\alpha/2}$ and $A^{\alpha/2}$. The finite constants cause no
problem: projecting $f_{i,n}$ off constants changes neither Gram matrix.

The subspace $A^{\alpha/2}c_{00}(G)$ is dense, since $A^{\alpha/2}$
is bounded and has dense range. Approximate any fixed number of top
eigenvectors of the compact operator $C$ by these images. Transport
the resulting Ritz problem and apply finite compression min--max.
Arbitrarily accurate approximation proves the assertion, including
eigenvalue multiplicities.
\end{proof}

\begin{theorem}[Local/global budget]\label{thm:budget}
Suppose the edges in Proposition~\ref{prop:ritz} are bridges, and
$\varepsilon_n\downarrow0$. With $\tau_{\mathrm w}$ the wired edge
resistance in the rooted limit,
\begin{equation}\label{eq:budget}
 \limsup_n\varepsilon_n\N_{C_n}(\varepsilon_n)
                \le\alpha\log\frac{2}{1+\tau_{\mathrm w}}.
\end{equation}
\end{theorem}
\begin{proof}
A finite unit bridge has leverage one. The determinant lemma on the
orthogonal complement of constants gives
\[
             \sum_j\log(1+\lambda_j(C_n))=\alpha\log2.
\]
For $q_n=\N_{C_n}(\varepsilon_n)$ and every fixed $k$, with zero
padding of eigenvalue lists,
\[
 (q_n-k)_+\log(1+\varepsilon_n)
    \le\alpha\log2-\sum_{j=1}^k\log(1+\lambda_j(C_n)).
\]
This also holds if $q_n<k$ since its left side is then zero.
Use Proposition~\ref{prop:ritz} and
$\varepsilon/\log(1+\varepsilon)\to1$, and then let $k\to\infty$.
Lemma~\ref{lem:fredholm} identifies the entire surviving local sum as
$\alpha\log(1+\tau_{\mathrm w})$.
\end{proof}

\begin{lemma}\label{lem:resistance}
If a rooted bridge limit has maximum degree $d\ge3$, then
$\tau_{\mathrm w}\ge2/d$.
\end{lemma}
\begin{proof}
Deleting the bridge leaves two halves, rooted at its endpoints, each
of root degree at most $d-1$. Write $r_+,r_-$ for their root-to-infinity
wired resistances, setting $r=\infty$ for a finite half.
Set $q=d-1$. In either infinite half, the cut between distance levels
$h,h+1$ has at most $q^{h+1}$ unit edges. The disjoint-cutset
Nash--Williams bound~\cite[Section 2.5]{LyonsPeres} gives
\[
                 r_\pm\ge\sum_{h=0}^\infty q^{-h-1}
                                =\frac1{d-2}.
\]
For completeness this also follows without a flow theorem. The
finitely supported radial function
\[
 f_R(x)=\frac{q^{-\operatorname{dist}(o,x)}-q^{-R}}{1-q^{-R}}
   \quad(\operatorname{dist}(o,x)<R),\qquad f_R(x)=0\text{ otherwise},
\]
equals one at the root and has energy at most
$(q-1)/(1-q^{-R})$. Thus the capacity is at most $d-2$.

The variational wired resistance of the bridge is
\[
 \tau_{\mathrm w}
 =\sup_{f\in c_{00}}\frac{|f(u)-f(v)|^2}{\langle f,Af\rangle}
 =\frac{r_++r_-}{1+r_++r_-},
\]
with value one if either half resistance is infinite. Indeed, at
fixed endpoint values $a,c$, the infimum energy is
$(a-c)^2+a^2/r_++c^2/r_-$. Minimize subject to $a-c=1$.
This proof also handles zero capacities by approximation with compactly
supported functions. The spectral variational formula identifies the
same supremum with $\|A^{-1/2}b\|^2$. The lower bounds on $r_\pm$
now give $\tau_{\mathrm w}\ge2/d$.
\end{proof}

\begin{proof}[Upper assertion of Theorem~\ref{thm:sharp}]
Choose any thresholds tending to zero and hosts maximizing the strict
count. The supremum is attained since the determinant bound makes the
integer count uniformly finite at each threshold. Bounded degree gives
an edge-rooted convergent subsequence. Its edited edge remains a bridge:
any alternative path in the limit would occur in a finite rooted ball.
A finite connected limit forces the hosts to be eventually that finite
graph, contributing zero to the scaled count. An infinite limit is
covered by Theorem~\ref{thm:budget} and Lemma~\ref{lem:resistance}, giving
$\alpha\log(2d/(d+2))$. For $d=3$ this is $\alpha\log(6/5)$.
\end{proof}

\section{Free/wired comparison and arbitrary existing edges}\label{sec:free-wired}
\begin{lemma}\label{lem:free-wired}
Let $G$ be an infinite connected simple unit-weight graph of maximum
degree $d\ge3$, and let $e=uv$ be any existing edge. If
$r=R^{\mathrm f}_G(u,v)$ and $w=R^{\mathrm w}_G(u,v)$ are its free
and wired resistances, then
\[
 w\ge\frac{r}{1+(d-2)r/2},\qquad
               \frac{1+r}{1+w}\le\frac{2d}{d+2}.
\]
\end{lemma}
\begin{proof}
Delete $e$ to obtain $H$, possibly disconnected. Each component of $H$
contains $u$ or $v$, and both roots have degree at most $d-1$.
For a finite induced exhaustion $H_R$ by distance at most $R$ from
$\{u,v\}$, let $S_R$ be the distance-$R$ shell. Define
\begin{align*}
 C_R^{\mathrm f}&=\inf\{E_{H_R}(f):f(u)-f(v)=1\},\\
 C_R^{\mathrm w}&=\inf\{E_{H_R}(f):f(u)-f(v)=1,\ f|_{S_R}=0\},\\
 p_R&=\inf\{E_{H_R}(f):f(u)=1,\ f|_{S_R}=0\},\\
 q_R&=\inf\{E_{H_R}(f):f(v)=1,\ f|_{S_R}=0\}.
\end{align*}
Free conductances increase to $C_H^{\mathrm f}$, the reciprocal of
the decreasing finite free resistances. The other three quantities
decrease to $C_H^{\mathrm w}$ and the singleton capacities $p,q$ of
$u,v$ in $H$. This follows directly by extending shell-zero functions
by zero outside $H_R$: their energies do not change, and their union
contains every finitely supported function. A fully exhausted finite
component causes no exception.

Eliminate every vertex of $H_R$ except $u,v$ and the \emph{individual}
vertices of $S_R$, before collapsing that shell. The elementary
star-elimination formula replaces conductances $c_i$ at a vertex by
nonnegative conductances $c_ic_j/(\sum_kc_k)$ between its neighbors.
Thus it preserves minimized energy and a nonnegative network.
Let $c$ be the reduced direct $uv$ conductance, and let $a,b$ be the
sums of conductances from $u,v$ to the shell. The free energy includes
the nonnegative term $c(f(u)-f(v))^2$, so $C_R^{\mathrm f}\ge c$.
With the shell grounded, its exact terminal energy is
$c(U-V)^2+aU^2+bV^2$. Minimizing the remaining root values gives
\[
 C_R^{\mathrm w}=c+\frac{ab}{a+b},\quad
 p_R=a+\frac{cb}{c+b}\ge a,\quad
 q_R=b+\frac{ca}{c+a}\ge b,
\]
where $0/0=0$. The function $xy/(x+y)$ is continuous and increasing
in both nonnegative arguments. Consequently
\[
 C_R^{\mathrm w}-C_R^{\mathrm f}
       \le\frac{ab}{a+b}\le\frac{p_Rq_R}{p_R+q_R}.
\]
Pass to the limit. The radial capacity trial in
Lemma~\ref{lem:resistance} applies separately at each root in $H$,
even when $H$ is connected: both roots have degree at most $d-1$.
It gives $p,q\le d-2$, so
\begin{equation}\label{eq:escape-conductance}
 C_H^{\mathrm w}-C_H^{\mathrm f}
             \le\frac{pq}{p+q}\le\frac{d-2}{2}.
\end{equation}
Restoring the unit edge adds one to both conductances. Thus, for
$k=(d-2)/2$, $1/w\le1/r+k$, proving the first inequality.
Here $0<r\le1$. The ratio therefore satisfies
\[
 \frac{1+r}{1+w}\le
 \frac{(1+r)(1+kr)}{1+(k+1)r}
     =1+\frac{kr^2}{1+(k+1)r}.
\]
The last expression is increasing for $0\le r\le1$, with value
$2d/(d+2)$ at one. This proves the second inequality.
\end{proof}

\begin{proof}[The all-edge upper bound]
For a finite host let $\tau_n=b_n^*L_n^\dagger b_n\le1$. Its exact
logarithmic budget is $\alpha\log(1+\tau_n)$. Pass to a rooted local
limit and a subsequence with $\tau_n\to\ell$. A finite local limit
again contributes zero to the scaled count. For an infinite limit $G$,
put $r=R_G^{\mathrm f}(u,v)$, $w=R_G^{\mathrm w}(u,v)$.
Every fixed finite connected subgraph $F\subset G$ containing $e$
embeds eventually into the finite hosts. Rayleigh monotonicity gives
$\tau_n\le R_F(u,v)$ eventually. Exhausting $G$ by such $F$ shows
$\ell\le r$. This is a free-resistance bound; replacing $r$ by $w$
would be incorrect.

Proposition~\ref{prop:ritz} and Lemma~\ref{lem:fredholm} do not require
bridges. Repeating the subtraction in Theorem~\ref{thm:budget} with the
finite budget $\alpha\log(1+\tau_n)$ gives
\[
 \limsup_n\varepsilon_n\N_{C_n}(\varepsilon_n)
 \le\alpha\log\frac{1+\ell}{1+w}
 \le\alpha\log\frac{1+r}{1+w}
 \le\alpha\log\frac{2d}{d+2}.
\]
Selecting maximizing finite hosts and extracting a rooted subsequence
proves the uniform all-edge upper. The bridge lower construction below
also belongs to the all-edge class, completing that part of
Theorem~\ref{thm:sharp}.
\end{proof}

\section{An exact finite target with equal reverse pivots}
The following separated-mass fact is the finite-dimensional ingredient
also used in the independent unweighted manuscript; its proof is
recalled to specify the uniformity required here.

\begin{lemma}[Separated masses]\label{lem:masses}
Let $K\succ0$, $K'=K+vv^*\succ0$ have size $m$. Put
\[
 D_\delta=\diag(\delta^{(m-1)/2},\ldots,1),\quad
 A_\delta=D_\delta K D_\delta,\quad B_\delta=D_\delta K'D_\delta.
\]
For $T_i=\{i,\ldots,m\}$ define
\[
 \kappa_i=\frac{\det K_{T_i}}{\det K_{T_{i+1}}},\quad
 \tau_i=v_{T_i}^*K_{T_i}^{-1}v_{T_i},\quad\tau_{m+1}=0.
\]
In an ordered eigenbasis of $A_\delta$, its relative correction tends
in norm to the diagonal entries
\[
 \left(\frac{\kappa_i'}{\kappa_i}\right)^\alpha-1
          =\left(\frac{1+\tau_i}{1+\tau_{i+1}}\right)^\alpha-1.
\]
The convergence is uniform over compact families of such SPD pairs.
\end{lemma}
\begin{proof}
Principal-minor expansion gives
$e_k(A_\delta)\sim\det K_{T_{m-k+1}}\delta^{k(k-1)/2}$.
To justify extracting eigenvalues from this expression, if $P_k$ is
the product of the largest $k$ eigenvalues, then
$P_k\le e_k\le\binom mkP_k$. Hence adjacent ordered eigenvalue
ratios are $O(\delta)$, using
$P_{k+1}P_{k-1}/P_k^2$ and the three elementary symmetric asymptotics.
All other products in $e_k$ are then negligible, and successive ratios
give $\lambda_i(A_\delta)\sim\kappa_i\delta^{m-i}$.

For $u_\delta=D_\delta v$, its leverage weights in that eigenbasis
are $z_i=|\langle u_\delta,\phi_i\rangle|^2/\lambda_i$.
Their sum is $\tau_1$. For every $a>0$,
\[
 u_\delta^*(A_\delta+\delta^{m-i}a)^{-1}u_\delta
   \longrightarrow v_{T_i}^*(K_{T_i}+ae_ie_i^*)^{-1}v_{T_i}
   =\tau_{i+1}+\frac{\kappa_i(\tau_i-\tau_{i+1})}{\kappa_i+a}.
\]
Earlier coordinates have divergent diagonal shifts and later ones
vanishing shifts. Comparing with the separated spectral expression
identifies $z_i\to\tau_i-\tau_{i+1}$.
In~\eqref{eq:resolvent}, scale $t=\lambda_i a$ for the $i$th diagonal.
Its denominator tends to
$1+\tau_{i+1}+(\tau_i-\tau_{i+1})/(1+a)$, and the envelope
$a^\alpha/(1+a)^2$ is integrable. The scalar fractional formula gives
the asserted diagonal limit. For $i<j$, dropping that denominator
bounds the off-diagonal entry in absolute value by
\[
 \sqrt{z_i z_j}\,x^{(1-\alpha)/2}\frac{1-x^\alpha}{1-x},
                     \qquad x=\lambda_i/\lambda_j\to0.
\]
It therefore vanishes. The determinant lemma on $K_{T_i}$ gives the
pivot-ratio expression. The same proof holds along convergent families
of SPD pairs. A compactness/subsequence argument proves uniformity.
\end{proof}

Fix the degree $d\ge3$, any integer $m\ge2$, and
$0<s<\log(2d/(d+2))$. Put
\[
 R_d=1+\frac2{d-2},\quad R'_d=\frac12+\frac2{d-2},\quad
 c=R_d^{-1},\quad c'={R'_d}^{-1},
\]
\[
 f(g)=\frac{c'+g}{c+g},\qquad r_i=\frac{2d}{d+2}e^{-is/m},\qquad
 g_i=\frac{c'-r_ic}{r_i-1}\quad(0\le i\le m).
\]
Here $g_0=0$, $g_i$ is strictly increasing, and every $g_i$ is finite
because $r_m>1$. Set $q_i=g_i-g_{i-1}>0$, $t_i=q_i^{-1}$,
$T=\diag(t_i)$, and define
\begin{equation}\label{eq:target}
 K=(T+R_d\one\one^*)^{-1},\qquad
 K'=(T+R'_d\one\one^*)^{-1}.
\end{equation}
These form a positive rank-one pair. The star identity
$K=\diag(q)-qq^*/(c+g_m)$ yields
\[
 \frac{\det K'_{T_i}}{\det K_{T_i}}
         =\frac{f(g_{i-1})}{f(g_m)},\qquad
 \log\frac{\kappa_i'}{\kappa_i}
         =\log f(g_{i-1})-\log f(g_i)=\frac{s}{m}.
\]
Lemma~\ref{lem:masses} therefore produces $m$ limiting relative
eigenvalues all equal to $e^{\alpha s/m}-1$.

\section{Realizing the target by bounded-degree unit-weight circuits}
\subsection{Repeated blocks and the exact impedance}
Choose any fixed rooted full binary tree $\mathcal B$ with exactly
$m$ labeled leaves; $m$ need not be a power of two. Define its PSD
current-energy matrix $R_{\mathcal B}$ by
\[
 j^*R_{\mathcal B}j
   =\sum_{a\in E(\mathcal B)}\left(\sum_{i\text{ below }a}j_i\right)^2.
\]
Let $q=d-1$ and $Q=q^h$, $h\ge1$. On the target side of the edited bridge, put
a complete depth-$h$ $q$-ary tree, identify each of its $Q$ leaves with
the root of a copy of $\mathcal B$, and connect every leaf labeled $i$
by a neck of $\ell_i$ unit edges to a distinct port of reservoir $i$.
All $Q$ ports with label $i$ belong to the same connected reservoir.
Choose integers $\ell_i\ge1$ with $\ell_i/Q\to t_i$.

On the other side put a complete depth-$h$ $q$-ary tree whose $Q$ leaves
are identified with distinct ports of a single anchor reservoir.
Join the two tree roots by the edited unit bridge. A port has internal
degree two and one external attachment. Top-tree degrees are at most $d$,
and bottom-block degrees are at most three,
and target leaves with necks have degree two. The whole host is simple,
connected and of maximum degree at most $d$; the edited edge is a bridge.

First give the reservoirs zero internal resistance and ground the
anchor. For terminal injections $j\in\R^m$, write $J=\sum_i j_i$.
The automorphisms permuting the repeated blocks transitively, and the
uniqueness of the minimum-energy flow, force each copy to carry
$j_i/Q$ at its leaf $i$. Each top tree costs
$J^2\sum_{k=1}^h q^{-k}=J^2(1-1/Q)/(d-2)$, the bottom
blocks cost $j^*R_{\mathcal B}j/Q$, and the necks cost
$\sum_i\ell_i j_i^2/Q$. Thus the exact grounded impedance, for bridge
resistance $r_b=1$ or $1/2$, is
\begin{equation}\label{eq:ideal-impedance}
 Z_Q(r_b)=\diag(\ell_i/Q)+\left(r_b+\frac{2(1-1/Q)}{d-2}\right)\one\one^*
                                      +Q^{-1}R_{\mathcal B}.
\end{equation}
This proves a matrix identity for arbitrary signed currents. It tends
to the two target impedances in~\eqref{eq:target}.

\subsection{Many ports, a uniform gap, and contact dilution}
By Bordenave's Theorem~1~\cite{Bordenave}, for every sufficiently large
even size there exists a simple cubic graph of Laplacian gap at least
some fixed $\gamma_0>0$. Subdivide $Q$ distinct edges once; their new
vertices are the ports. These modified reservoirs still have gap
\begin{equation}\label{eq:reservoir-gap}
                    \gamma=(2+14/\gamma_0)^{-1}>0
\end{equation}
independent of $Q$ and the reservoir volume. Indeed the old energy is
at most twice the new energy $E$, so old-vertex variance about the old
mean is at most $2E/\gamma_0$. Assign each port to one old endpoint;
at most three ports are assigned to any old vertex. Their total squared
deviation from that mean is at most twice $E$ plus six times the old
variance. Adding old and new vertices proves~\eqref{eq:reservoir-gap}.
Variance about the new mean is smaller. For fixed $Q$, the available
sizes (old even size plus $Q$) approximate all positive prescribed
mass ratios after common scaling.

Let $P_i$ select the $Q$ distinct ports of reservoir $i$, and define
its contact matrix
\[
 H_i=P_i^*L_{\mathrm{res},i}^\dagger P_i,
                     \qquad0\preceq H_i\preceq R I,\quad R=1/\gamma.
\]
If its port current vector is $a$, the internal load is
$P_ia-(\one^*a)\one/n_i$, since a terminal current in the mean
coordinate is supplied uniformly. The minimum internal flow energy
is exactly $a^*H_i a$.

Let $K_Q(H;r_b)$ be the true grounded stiffness on target reservoir
means, with anchor mean zero. Dropping the nonnegative reservoir
energies, and then using the ideal symmetric flow as a trial flow in
the true network, gives the uniform matrix inequalities
\begin{equation}\label{eq:contact-bound}
 0\preceq K_Q(H;r_b)^{-1}-Z_Q(r_b)
             \preceq\frac RQ(I+\one\one^*).
\end{equation}
The trial currents are $j_i/Q$ at every target-$i$ port and $-J/Q$
at every anchor port. Their total contact cost is at most
$(R/Q)(\sum_i j_i^2+J^2)$. This proves~\eqref{eq:contact-bound}
uniformly over all possible contact matrices, without assuming any
contacts converge. It also justifies the use of means rather than
one particular physical port as the terminals.

Consequently, uniformly over the compact contact set
$0\preceq H_i\preceq RI$,
\[
 K_Q(H;1)\longrightarrow K,\qquad K_Q(H;1/2)\longrightarrow K'.
\]
In fact the true impedance satisfies the exact identity
$K_Q(H;1/2)^{-1}=K_Q(H;1)^{-1}-\tfrac12\one\one^*$:
every target current vector forces bridge current $J$, and all other
minimum-flow terms are independent of bridge resistance. Inverting
therefore gives a positive rank-one change in terminal stiffness.
The limiting pair is SPD, and
all reverse determinant ratios are continuous nearby. For any
$0<\eta<s$, choose one finite $Q$ so that every allowed contact tuple
satisfies
\begin{equation}\label{eq:almost-pivots}
 \log(\kappa_i'(H)/\kappa_i(H))\ge\frac{s-\eta}{m}
                              \qquad(1\le i\le m).
\end{equation}

\section{From terminal masses to ordinary unweighted graphs}
For this fixed finite core, give target $i$ mass
$M_i=\delta^{-(m-i)}$. Lemma~\ref{lem:masses}, uniformly on the compact
contact family, makes all $m$ grounded relative eigenvalues arbitrarily
close from below to at least $e^{\alpha(s-\eta)/m}-1$ when $\delta$
is sufficiently small. Give the anchor mass $M_0$ and let it tend to
infinity. The full mass-normalized $(m+1)$-terminal pair converges to
zero direct-sum the grounded pair. Its one common kernel converges to
the anchor coordinate; its other eigenvalues stay positive. Finite
functional calculus, including inverse powers on support, transfers
the $m$ lower eigenvalues for sufficiently large finite $M_0$, uniformly
in the same contact family.

The final passage to unit vertex masses needs a fractional, rather
than merely a static, reservoir limit.
\begin{lemma}[Fixed-core multiport reservoir realization]\label{lem:reservoir}
Fix a finite core, finitely many ports per reservoir, uniform internal
reservoir gap, and positive prescribed masses $M_0,\ldots,M_m$.
Take reservoir sizes $n_i(a)/a\to M_i$, and let $P_a$ be the isometry
of normalized reservoir indicators. Along a contact-convergent
subsequence, let $H,H'$ be the prescribed-mass static Schur pair. Then
\[
 P_a^*(aL_a)^\alpha P_a\longrightarrow H^\alpha,\qquad
 P_a^*(aL'_a)^\alpha P_a\longrightarrow {H'}^\alpha.
\]
Every eigenvalue of the relative prescribed-mass pair is a limit of a
generalized Ritz eigenvalue of the full graph response and hence a
lower bound, up to vanishing error, for the full correction eigenvalues.
\end{lemma}
\begin{proof}
Decompose into the indicator span and its orthogonal complement, which
contains the mean-zero reservoir directions and all core coordinates:
\[
 L_a=\begin{pmatrix}A_a&B_a\\B_a^*&F_a\end{pmatrix},\qquad
 \|A_a\|=O(a^{-1}),\quad\|B_a\|=O(a^{-1/2}),\quad F_a\succeq\xi I.
\]
Here $\xi>0$ is independent of reservoir sizes. Internal gaps control
mean-zero reservoir norms and their fixed number of port deviations.
The fixed core's positive Dirichlet stiffness controls its norm by
actual external energy plus those deviations. Combining these bounds
proves the last inequality, also for the updated blocks.

The static Schur matrix $S_a=A_a-B_aF_a^{-1}B_a^*$ is exactly the true
stiffness on reservoir means conjugated by $\diag(n_i^{-1/2})$.
Thus $aS_a\to H$ and $aS'_a\to H'$ along the given subsequence.
For each fixed $t>0$, block inversion yields
\begin{align*}
 G_a(t)&=P_a^*(aL_a+t)^{-1}P_a\\
 &=\left[aS_a+t+
   aB_a\{F_a^{-1}-(F_a+t/a)^{-1}\}B_a^*\right]^{-1}
       \longrightarrow(H+t)^{-1}.
\end{align*}
The correction inside the brackets is $O(t/a)$. Moreover
$P_a^*aL_aP_a=aA_a$ is uniformly bounded. Spectral calculus therefore
gives the uniform positive bounds
\[
 0\preceq I-tG_a(t)
   =P_a^*aL_a(aL_a+t)^{-1}P_a
   \preceq I,\qquad
 I-tG_a(t)\preceq \frac Ct I.
\]
The fractional formula
\[
 P_a^*(aL_a)^\alpha P_a
       =c_\alpha\int_0^\infty t^{\alpha-1}(I-tG_a(t))\,dt
\]
now passes to the limit by the integrable envelope
$t^{\alpha-1}\min(1,C/t)$, since $0<\alpha<1$. This proves both
fractional compression limits.

The common kernel of $H,H'$ is the vector
$z=(\sqrt{M_0},\ldots,\sqrt{M_m})^*$.
Choose an isometry $V:\R^m\to z^\perp$. The generalized Gram matrices
on the graph test vectors $P_aV$, multiplied by the common factor
$a^\alpha$, converge to
\[
 E=V^*H^\alpha V\succ0,\qquad
 D=V^*({H'}^\alpha-H^\alpha)V.
\]
Hence their generalized eigenvalues converge to those of the
prescribed-mass relative pair. As in Proposition~\ref{prop:ritz},
finite constants can be projected away without changing either Gram
matrix, and compression min--max transfers the lower eigenvalues.
No convergence of inverse powers or high fractional response is needed.
\end{proof}

Contacts lie in a fixed finite-dimensional compact set, since $Q$ was
fixed before the volume limit. If approximation to the contact-dependent
mass pair failed uniformly, a contact-convergent subsequence would
contradict Lemma~\ref{lem:reservoir}. Thus one sufficiently large finite
volume scale realizes all $m$ lower eigenvalues in the full graph.

\begin{proof}[Lower assertion of Theorem~\ref{thm:sharp}]
Fix $s<\log(2d/(d+2))$ and $0<\eta<s$. For every integer $m\ge2$, the finite
choices above produce graph eigenvalues as close from below as desired
to $e^{\alpha(s-\eta)/m}-1$. Given small $\varepsilon$, choose an integer
\[
 m<\frac{\alpha(s-\eta)}{\log(1+\varepsilon)},\qquad
 m\log(1+\varepsilon)\longrightarrow\alpha(s-\eta).
\]
The target eigenvalue has a strictly positive margin above
$\varepsilon$. Choose successively the finite $Q$, mass separation,
anchor mass, and reservoir sizes so that all errors fit that margin.
Then $W^{\mathrm{br}}_{d,\alpha}(\varepsilon)\ge m$. Therefore
\[
 \liminf_{\varepsilon\downarrow0}\varepsilon
            W^{\mathrm{br}}_{d,\alpha}(\varepsilon)
                         \ge\alpha(s-\eta).
\]
Let $\eta\downarrow0$ and $s\uparrow\log(2d/(d+2))$, and combine with the
proved upper bound.
\end{proof}

\subsection{Bipartite hosts}
To obtain the same coefficient for bipartite hosts, use bipartite
cubic reservoirs. Bordenave's two-sided adjacency bound supplies
simple cubic graphs with all nontrivial adjacency eigenvalues in
$[-3+\gamma_0,3-\gamma_0]$. Their bipartite double covers are simple,
connected cubic graphs with the same positive Laplacian gap.

Replace each of $Q$ distinct cover edges by a path of length three.
The two new internal vertices have opposite colors; select one as the
port. These reservoirs retain a uniform gap
$\gamma_{\mathrm{bip}}=(2+21/\gamma_0)^{-1}$: old energy is at most
$3E$, old variance at most $3E/\gamma_0$, and assigning each new
vertex to its adjacent old endpoint assigns at most three new vertices
to each old vertex. New-vertex variance about the old mean is at most
$2E+6V_{\mathrm{old}}$. Adding the old variance gives the asserted
constant. Thus $0\preceq H_i\preceq\gamma_{\mathrm{bip}}^{-1}I$
uniformly. The available sizes are $4k+2Q$, which approximate any
fixed positive mass ratios after scaling, and there are enough distinct
edges once $k$ is large.

The exterior core is a tree. All anchor leaves have the same depth,
and each fixed bottom-block label has the same depth in every repeated
copy. Thus every attachment of one reservoir requires the same port
color. On each subdivided edge choose the internal vertex of that
color. This realizes a consistent bipartition, including all necks.
Alternatively one may adjust each $\ell_i$ by at most one to prescribe
its parity, without changing $\ell_i/Q\to t_i$. Every port still has
internal degree two and one attachment. The same contact bound, mass
choices and realization lemma prove the lower coefficient in the
bipartite subclass; the upper bound already applies to it.

\section{Interpretation and verification status}
The cubic-tree local response has wired leverage $2/3$, and its
entire determinant is $(5/3)^\alpha$. Every finite edited bridge has
determinant $2^\alpha$. The ratio $(6/5)^\alpha$ is exactly the budget
that the construction can distribute over many vanishing global modes.
The repeated ports, rather than separately ordered individual leaves,
are essential to the lower proof: one mass coordinate aggregates $Q$
small leaf contributions. No lower bound on the last single-leaf
admittance increment is assumed.

The analytic ingredients are standard resolvent functional calculus,
trace-class determinant calculus and min--max. The sole external graph
existence input is the cubic expander statement of Bordenave's
Theorem~1, whose specialization is stated explicitly above. The
separated-mass and reservoir-limit arguments are included here.
Independent session notes separately record the operator proof,
rooted upper proof, and lower construction. A numerical check script
verifies the repeated-block impedance and the tuned pivots on small
finite circuits; its checks supplement, rather than prove, the argument.
This manuscript is a session research result awaiting the root task's
additional independent audits; no historical-priority claim is made.

\begin{thebibliography}{9}
\bibitem{Bordenave} C.~Bordenave. A new proof of Friedman's second
eigenvalue theorem and its extension to random lifts. Ann. Sci. Ecole
Norm. Sup. 53 (2020), 1393--1439. Theorem~1.
\url{https://arxiv.org/abs/1502.04482}.
\bibitem{Bornemann} F.~Bornemann. On the numerical evaluation of Fredholm
determinants. Math. Comp. 79 (2010), 871--915. Sections~3--4.
\url{https://arxiv.org/abs/0804.2543}.
\bibitem{LyonsPeres} R.~Lyons and Y.~Peres. Probability on Trees and
Networks. Cambridge University Press, 2016. Section~2.5 and Chapter~9.
Author-hosted editions: \url{https://rdlyons.pages.iu.edu/prbtree/}.
\end{thebibliography}
\end{document}
